\section{Approche réflexive}
\subsection{Les connaissances/compétences acquises lors de la rédaction de cet incident clinique ;}
\begin{itemize}[label=\textbullet]
    \item Approfondir ma capacité de recherche dans la littérature car celle-ci était pauvre pour cette thématique. Rechercher de nouvelles bases de données et aiguiser d’autant plus mon esprit critique.
    \item La déconstruction sur certaines croyances vis-à-vis du déni de grossesse. J’ai eu l’occasion de développer la physiopathologie et d’améliorer la façon dont je souhaitais aborder les patientes expérimentant un déni. 
    \item Ce cas clinique a montré l’importance de la pluridisciplinarité. Au-delà de l’équipe soignante, il y avait une collaboration étroite avec l’assistante sociale et avec la Juge responsable. L’aspect social étant vraiment présent et imbriqué dans la santé et la clinique du patient que l’assistant social devient un membre primordial et fait partie intégrante de la prise en charge. 
    \item Observer de mes yeux l’absence de lien mère-enfant a provoqué en moi beaucoup d’émotions et d’incompréhension. Il m’a fallu apprendre à prendre du recul sur la situation et garder un professionnalisme sans laisser transparaître mes sentiments et moi-même me protéger de cet impact.
\end{itemize}

\subsection{Les influences probables de ce travail sur sa pratique médicale future ;}

\begin{itemize}[label=\textbullet]
    \item Ne pas stigmatiser les patients selon leur milieu social et prendre le temps de comprendre la problématique et le contexte. Celui-ci peut aider à la prise en charge et savoir comment communiquer et établir une relation de confiance avec le patient et sa famille.
    \item Préserver la curiosité d’en apprendre plus sur les pathologies que l’on rencontre en clinique. Se remettre en question sur les bases acquises et rester à jour au niveau du savoir médical.
\end{itemize}

\subsection{Les difficultés (dépendantes ou indépendantes de sa volonté) ayant entravé son acquisition de nouvelles connaissances/compétences au cours de cette expérience ;}

\begin{itemize}[label=\textbullet]
    \item La restriction de l’information contenue dans les protocoles ou rapport présent dans le dossier. Les anamnèses d’entrée sont clairement centrées sur la clinique, toutefois, vu le contexte, avoir davantage d’informations sur l’aspect social aurait été utile.
    \item En lien avec l’émotion ressentie, j’ai également ressenti beaucoup d’impuissance. Je pense qu’il est naturel d’avoir un réflexe correcteur face à ce patient si jeune qui subit beaucoup face à un environnement nuisible. Je peux ajouter que face au déni des parents, il était parfois difficile de trouver mes mots afin de rester professionnel et d’avoir les termes appropriés et juste dans le sens que e-je voulais leur donner.
\end{itemize}

\subsection{Les solutions que l’étudiant·e peut mettre en place pour remédier à ces difficultés.}

\begin{itemize}[label=\textbullet]
    \item Pour améliorer les rapports, je peux penser à ajouter de façon systématique un contexte social et détailler les circonstances d’admission. 
    \item Il est important de déconstruire ce réflexe qui parfois entrave la bonne prise en charge et la relation thérapeutique. Le premier pas est d’abord de le reconnaitre et d’en prendre conscience. Ensuite, si le besoin en est, communiquer de cette problématique à un collègue voire à un professionnel afin d’être accompagné dans cette démarche. Plusieurs ouvrages existent et aident le soignant qui prend la peine de les lire. 
    \item D’un point de vue purement communicationnel, je pense qu’il me reste beaucoup à apprendre et la formation peut-être multimodal via des cours, des ouvrages, des conférences,…
    \end{itemize}


